% concatenated from Kolmogorov-Lewko-2.tex
\documentclass[11pt]{article}
\usepackage[T1]{fontenc}
\usepackage{lmodern}
\usepackage[margin=1in]{geometry}
\usepackage{amsmath,amssymb,amsthm,mathtools,microtype}
\usepackage{needspace}
\usepackage{xcolor}
\usepackage[colorlinks=true,linkcolor=blue!45!black,citecolor=blue!45!black,urlcolor=blue!45!black]{hyperref}
\hypersetup{pdftitle={On Kolmogorov's rearrangement problem and Garsia's conjecture},pdfauthor={Mark Lewko}}
\newtheorem{theorem}{Theorem}
\newtheorem{proposition}[theorem]{Proposition}
\newtheorem{lemma}[theorem]{Lemma}
\newtheorem{corollary}[theorem]{Corollary}
\theoremstyle{remark}
\newtheorem{remark}[theorem]{Remark}
\numberwithin{equation}{section}
\newcommand{\T}{\mathbb T}
\newcommand{\N}{\mathbb N}
\newcommand{\Z}{\mathbb Z}
\newcommand{\R}{\mathbb R}
\newcommand{\C}{\mathbb C}
\newcommand{\E}{\mathbb E}
\newcommand{\PP}{\mathbb P}
\newcommand{\one}{\mathbf 1}
\newcommand{\calM}{\mathcal M}
\newcommand{\calF}{\mathcal F}
\newcommand{\calA}{\mathcal A}
\newcommand{\calB}{\mathcal B}
\newcommand{\calE}{\mathcal E}
\newcommand{\norm}[1]{\left\lVert#1\right\rVert}
\setlength{\emergencystretch}{2em}
\allowdisplaybreaks[2]
\title{On Kolmogorov's rearrangement problem\\ and Garsia's conjecture}
\author{Mark Lewko}
\date{}
\begin{document}
\maketitle

\begin{abstract}
We give negative answers to Kolmogorov's rearrangement problem and Garsia's conjecture. We construct a complete uniformly bounded orthonormal system for which every rearrangement admits a square-summable series divergent almost everywhere. The construction is built from two copies of the trigonometric system in different orderings. The main ingredient is a combinatorial lemma which finds a prescribed permutation pattern as a subsequence of at least one of two longer permutations. Its proof uses Szemer\'edi's theorem and a counting argument.
\end{abstract}

\section{Introduction}

Let $(\Omega,\mu)$ denote a probability space and let $\{\phi_n\}_{n\geq1}$ be an orthonormal system (ONS) on $\Omega$. For $(a_n)\in\ell^2$, orthogonality guarantees $L^2$ convergence of the series $\sum_n a_n\phi_n$. We call an ONS a \emph{system of convergence} if every square-summable series converges almost everywhere. Kolmogorov's rearrangement problem asks whether every ONS can be permuted to become a system of convergence.\footnote{The problem is generally attributed to Kolmogorov. Bourgain \cite{Bourgain} cites Kolmogorov's 1924 paper \cite{Kolmogorov1924}, but that paper concerns lacunary Fourier convergence and does not state the rearrangement problem.} For general background on orthogonal series, see Olevskii's book \cite{OlevskiiBook}.

The trigonometric system provides a natural example. Kolmogorov stated without proof, as Theorem~III of his 1927 paper with Menshov \cite{KolmogorovMenshov}, that an $L^2$ Fourier series can have an almost everywhere divergent rearrangement. Zahorski published a proof sketch in 1960 \cite{Zahorski}; see also \cite{BKL,Karagulyan}. In 1964, Olevskii \cite{Olevskii1964} constructed the first complete uniformly bounded ONS that is a system of convergence. In 1966, Carleson \cite{Carleson} proved that the trigonometric system in its natural ordering is a system of convergence. Thus permuting the trigonometric system can change whether it is a system of convergence.

For a general ONS, Garsia proved that each square-summable coefficient sequence admits an almost everywhere convergent rearrangement \cite{Garsia1964,Garsia1970}. Here the permutation may depend on the coefficients. Kolmogorov's problem asks for one permutation which works for all square-summable coefficient sequences.

Garsia's conjecture \cite{Garsia1970} is a stronger finite formulation. It asks whether every finite ONS $\{\phi_n\}_{n=1}^N$ admits a permutation $\tau$ of $\{1,\ldots,N\}$ such that
\begin{equation}\label{eq:garsia}
 \left\|\max_{q\leq N}\left|\sum_{j\leq q}
 a_{\tau(j)}\phi_{\tau(j)}\right|\right\|_2
 \leq C\left(\sum_{n=1}^N|a_n|^2\right)^{1/2}
\end{equation}
for all coefficients $\{a_n\}$, with $C$ independent of the system and $N$. This conjecture would imply a positive answer to Kolmogorov's problem.

In 1989, Bourgain \cite{Bourgain} proved that if $|\phi_n|\leq A$, then some permutation satisfies \eqref{eq:garsia} with $C=C_A\log\log(N+3)$, where $C_A$ depends only on $A$. His proof uses random rearrangements and estimates for stochastic processes, including entropy and chaining arguments related to his work on the $\Lambda(p)$ problem \cite{BourgainLambda}. Bourgain also showed, by a combinatorial argument \cite[Section~2]{Bourgain}, that a positive answer to Kolmogorov's problem would imply almost everywhere convergence of the usual Walsh partial sums of $L^2$ functions. This is the Walsh analogue of Carleson's theorem, proved by Billard in 1967 \cite{Billard}. Our work here is inspired by the combinatorial ideas underlying this observation of Bourgain.

Our main result is the following resolution of Kolmogorov's problem and Garsia's conjecture:

\Needspace{5\baselineskip}
\begin{theorem}\label{thm:main}
\leavevmode
\begin{enumerate}
\item There is an absolute constant $c>0$ with the following property. For every $H>0$, there is a finite ONS $\{\phi_n\}_{n=1}^N$ on a probability space $(\Omega,\mu)$, with $|\phi_n|=1$, such that for every permutation $\tau$ of $\{1,\ldots,N\}$ there are real coefficients with $\sum_{n=1}^N a_n^2=1$ satisfying
\[
 \mu\left\{x\in\Omega:\max_{q\leq N}
 \left|\sum_{j\leq q}a_{\tau(j)}\phi_{\tau(j)}(x)\right|>H\right\}\geq c.
\]
\item There is a complete ONS $\{\psi_n\}_{n\geq1}$ on $[0,1]$, with $|\psi_n|=1$, such that for every permutation $\tau$ of $\N$ there are real coefficients $(a_n)\in\ell^2$ for which
\[
 \sup_{q\geq1}\left|\sum_{j=1}^q a_{\tau(j)}\psi_{\tau(j)}(x)\right|=\infty
\]
for almost every $x\in[0,1]$.
\end{enumerate}
There are also real-valued examples for both assertions, uniformly bounded by $\sqrt2$.
\end{theorem}

The first assertion disproves Garsia's conjecture: after any permutation the maximal operator has norm at least $H\sqrt c$. The second gives one infinite system for which every permutation admits a square-summable series divergent almost everywhere, and hence answers Kolmogorov's problem negatively.

The proof as written gives no useful bound for the dependence of $N$ on $H$, since it uses only qualitative forms of the Fourier divergence theorem and Szemer\'edi's theorem. The finite examples may be enlarged to any greater size by extending $\sigma$ to fix the additional frequencies. Together with Bourgain's bound, the theorem therefore shows that the optimal universal constant in \eqref{eq:garsia} for uniformly bounded systems of size $N$ tends to infinity with $N$, but no faster than $\log\log N$. We do not know its true rate of growth.

The finite construction uses two copies of the trigonometric system. Let $\T:=\R/\Z$ denote the unit circle with normalized Lebesgue measure, and write $e(t):=e^{2\pi it}$ and $[N]:=\{1,\ldots,N\}$. Given a permutation $\sigma$ of $[N]$, set
\begin{equation}\label{eq:construction}
 \phi_n(x,0):=e(nx),\qquad
 \phi_n(x,1):=e(\sigma(n)x),\qquad n\in[N],
\end{equation}
on $\T\times\{0,1\}$, where each copy of $\T$ has measure $1/2$. These functions are orthonormal and have absolute value one. We choose $\sigma$ using the combinatorial lemma in the next section, then apply the classical divergence theorem for rearranged Fourier series.

\section{The combinatorial argument}

Translating the frequencies of a trigonometric polynomial by an integer, or dilating them by a positive integer, preserves the distribution of its maximal absolute partial sum. A finite Fourier example can therefore be placed on any arithmetic progression, provided that the permutation of its frequencies is preserved. This is why arithmetic progressions enter the construction.

Fix a positive integer $m$ and a permutation $\pi$ of $[m]$. In a permuted list of integers, we seek an arithmetic progression $x_1<\cdots<x_m$ whose terms occur as $x_{\pi(1)},\ldots,x_{\pi(m)}$. Other integers in the list may occur between these terms. The following lemma provides the permutation needed for the two-copy construction.

\begin{lemma}\label{lem:two}
Let $\pi$ be a permutation of $[m]$. There exist $N\geq m$ and a permutation $\sigma$ of $[N]$ such that, for every permutation $\tau$ of $[N]$, there is an $m$-term arithmetic progression $x_1<\cdots<x_m$ in $[N]$ for which
\[
 x_{\pi(1)},\ldots,x_{\pi(m)}
\]
is a subsequence of at least one of
\[
 \tau(1),\ldots,\tau(N)
 \qquad\text{and}\qquad
 \sigma(\tau(1)),\ldots,\sigma(\tau(N)).
\]
\end{lemma}

For the system \eqref{eq:construction}, a permutation $\tau$ of the functions gives the frequency permutations $\tau$ and $\sigma\circ\tau$ on the two copies of $\T$.

We prove the lemma in two steps: first count the colorings that fail to contain the required progression, then use this count to choose $\sigma$. Here $m$ remains the progression length. We introduce $K\geq m$ available colors, labeled $1,\ldots,K$; the chosen progression may use any $m$ distinct colors. The counting estimate holds for every fixed $K\geq m$, and the final step will use $K=m^2$.

\begin{lemma}\label{lem:count}
Fix integers $m\geq2$ and $K\geq m$, and a permutation $\pi$ of $[m]$. All but at most
\begin{equation}\label{eq:count}
 (m-1)^N\exp(o(N))
\end{equation}
of the $K^N$ colorings $c:[N]\to[K]$ admit an $m$-term arithmetic progression $x_1<\cdots<x_m$ in $[N]$ with distinct colors such that $x_{\pi(j)}$ has the $j$-th smallest color among these $m$ points, for each $j\in[m]$. Equivalently,
\begin{equation}\label{eq:pattern}
 c(x_{\pi(1)})<\cdots<c(x_{\pi(m)}).
\end{equation}
Here $N\to\infty$ with $m$, $K$, and $\pi$ fixed.
\end{lemma}

For the application to Lemma~\ref{lem:two}, where $K=m^2$, the point is that the square of this bound is exponentially smaller than $K^N=m^{2N}$. The bound is sharp up to the $\exp(o(N))$ factor, since a coloring using at most $m-1$ colors admits no progression with $m$ distinct colors.

\begin{proof}
Let $\calF_N$ be the family of colorings $c:[N]\to[K]$ containing no progression satisfying \eqref{eq:pattern}. We count these colorings by deleting the points $N,N-1,\ldots,1$ in turn.

\emph{The counting step.}
Consider any family $\calA$ of colorings of $[r]$, where $1\leq r\leq N$. We delete the last point $r$. For each color $i\in[K]$, let
\[
 \calA_i:=\{c|_{[r-1]}:c\in\calA,\ c(r)=i\}.
\]
Here $[0]:=\varnothing$. Thus $\calA_i$ consists of the colorings of $[r-1]$ that extend to a member of $\calA$ by giving $r$ color $i$. We claim that
\begin{equation}\label{eq:recursion}
 |\calA|\leq(m-1)\left|\bigcup_{i=1}^K\calA_i\right|
 +\sum_{\substack{Q\subseteq[K]\\|Q|=m}}
 \left|\bigcap_{i\in Q}\calA_i\right|.
\end{equation}
Indeed, a coloring of $[r-1]$ that belongs to exactly $s$ of the families $\calA_i$ has $s$ extensions in $\calA$. If $s\geq1$, it contributes $m-1+\binom sm$ on the right, and this is at least $s$: for $s\leq m$ this is clear, and for $s>m$ we have $\binom sm\geq\binom s1=s$. If $s=0$, it contributes zero to both sides.

Apply \eqref{eq:recursion} first with $\calA=\calF_N$ and $r=N$, then to each family on the right with $r=N-1$, and continue. At each application, $\calA$ denotes the family currently being counted, and $\calA_i$ is defined from that family by deleting its last point. A union contributes a factor $m-1$; an intersection contributes a choice of an $m$-element color set $Q$.

\Needspace{6\baselineskip}
\emph{Controlling the intersections.}
Fix one term in the resulting sum. Let $S\subseteq[N]$ be the points at which an intersection was chosen, and let $Q_x$ be the color set chosen at $x\in S$. These choices determine the term. Once all points have been deleted, the remaining family is either empty or consists of the unique coloring of the empty set. A nonzero term therefore has weight $(m-1)^{N-|S|}$.

For such a term, we may prescribe any color from $Q_x$ at each $x\in S$ and extend these choices to a coloring in $\calF_N$. To see this, restore the points in reverse order. Membership in a union allows at least one color at the restored point; membership in the chosen intersection allows every color in $Q_x$, so we use the prescribed one. The colors outside $S$ may depend on these prescriptions.

Fix an $m$-element color set $Q=\{q_1<\cdots<q_m\}$. The points $x\in S$ with $Q_x=Q$ contain no $m$-term arithmetic progression. Otherwise, on such a progression $x_1<\cdots<x_m$, prescribe color $q_j$ at $x_{\pi(j)}$. The extension just described would belong to $\calF_N$ and satisfy \eqref{eq:pattern}, a contradiction.

Let $r_m(N)$ denote the largest size of a subset of $[N]$ containing no $m$-term arithmetic progression. By Szemer\'edi's theorem \cite{Szemeredi}, $r_m(N)=o(N)$. Each of the $\binom Km$ classes of points with the same color set has size at most $r_m(N)$, and hence
\[
 |S|\leq\binom Km r_m(N)=o(N)
\]
uniformly over the nonzero terms. Let $d_N$ be the largest possible $|S|$ among these terms; thus $d_N=o(N)$.

For $|S|=j$, there are $\binom Nj$ choices of $S$ and $\binom{K}{m}^{j}$ choices of its color sets. Summing the weights gives
\[
 |\calF_N|\leq\sum_{j=0}^{d_N}\binom Nj\binom{K}{m}^{j}(m-1)^{N-j}
 \leq(m-1)^N\exp(o(N)),
\]
where the last estimate uses $d_N=o(N)$.
\end{proof}

\begin{proof}[Proof of Lemma~\ref{lem:two}]
For $m=1$, take $N=1$. For $m\geq2$, set $K:=m^2$ and take $N$ divisible by $K$. Call a coloring $c:[N]\to[K]$ \emph{balanced} if each color is used exactly $N/K$ times, and let $\calB_N$ be the family of balanced colorings. Then
\[
 |\calB_N|=:T_N=\frac{N!}{((N/K)!)^K}=K^N\exp(o(N)).
\]
Let $\calF_N$ be the exceptional family in Lemma~\ref{lem:count}. For a fixed balanced coloring $c$ and a uniformly random permutation $\sigma$ of $[N]$, the coloring $c\circ\sigma^{-1}$ is uniform on $\calB_N$. Thus the expected number of balanced colorings $c$ for which both $c$ and $c\circ\sigma^{-1}$ belong to $\calF_N$ is
\[
 \frac{|\calF_N\cap\calB_N|^2}{T_N}
 \leq\frac{|\calF_N|^2}{T_N}
 \leq\left(\frac{(m-1)^2}{m^2}\right)^N\exp(o(N))\longrightarrow0.
\]
For $N$ sufficiently large, choose $\sigma$ for which there is no such coloring.

Now fix an arbitrary permutation $\tau$ of $[N]$. Partition the positions in the sequence $\tau(1),\ldots,\tau(N)$ into $K$ consecutive intervals, each of length $N/K$. Define the coloring $c_\tau:[N]\to[K]$ by
\begin{equation}\label{eq:block-coloring}
 c_\tau(\tau(j))=\ell
 \quad\text{when}\quad
 (\ell-1)\frac{N}{K}<j\leq\ell\frac{N}{K},
 \qquad 1\leq\ell\leq K.
\end{equation}
Thus the color of an integer records the interval in which it appears in this particular permutation. Define $c_{\sigma\circ\tau}$ by the same rule for the sequence $\sigma(\tau(1)),\ldots,\sigma(\tau(N))$. These colorings are balanced, and
\[
 c_{\sigma\circ\tau}(\sigma(n))=c_\tau(n),
 \qquad n\in[N].
\]
In particular, $c_{\sigma\circ\tau}=c_\tau\circ\sigma^{-1}$. By the choice of $\sigma$, at least one of $c_\tau$ and $c_{\sigma\circ\tau}$ lies outside $\calF_N$. Lemma~\ref{lem:count} supplies a progression whose terms, permuted by $\pi$, have strictly increasing colors for that coloring. By \eqref{eq:block-coloring}, smaller colors correspond to earlier positions in the corresponding permutation. Hence $x_{\pi(1)},\ldots,x_{\pi(m)}$ is the required subsequence.
\end{proof}

\section{Proof of Theorem~\ref{thm:main}}

We first record a finite consequence of the classical divergence theorem for rearranged Fourier series \cite{KolmogorovMenshov,Zahorski}.

\begin{lemma}\label{lem:fourier}
There is $\delta>0$ such that, for every $H>0$, there are $m$, a permutation $\pi$ of $[m]$, and real coefficients $b_1,\ldots,b_m$ satisfying
\begin{equation}\label{eq:fourier}
 \sum_{j=1}^m b_j^2=1,\qquad
 \left|\left\{x\in\T:\max_{q\leq m}
 \left|\sum_{j\leq q}b_j e(\pi(j)x)\right|>H\right\}\right|\geq\delta.
\end{equation}
\end{lemma}

\begin{proof}
Fix a divergent rearranged $L^2$ Fourier series, written using the characters $e(kx)$, $k\in\Z$, and let $T_q$ denote its partial sums. For $s\geq1$, let $E_s$ be the set of points $x$ such that, for every $r$, some $v>u\geq r$ satisfy
\[
 |T_v(x)-T_u(x)|>\frac1s.
\]
The union of the sets $E_s$ has full measure. Choose $s$ such that $|E_s|=:\eta>0$, and set $\varepsilon:=1/s$. For every $r$,
\[
 \sup_{q>r}|T_q-T_r|>\frac{\varepsilon}{2}
\]
on $E_s$. Choose $r$ so that the $\ell^2$ norm of the coefficient tail is smaller than $\varepsilon/(4H)$. The sets
\[
 \left\{x:\max_{r<q\leq M}|T_q(x)-T_r(x)|>\frac{\varepsilon}{2}\right\},
 \qquad M>r,
\]
increase to a set containing $E_s$. Hence some $M$ gives a set of measure at least $\eta/2$. Normalize the coefficient vector from the finite tail $T_M-T_r$. Its maximal partial sum exceeds $2H$ on this set.

Write this normalized coefficient vector as $u+iv$, where $u$ and $v$ are real vectors. By the triangle inequality, at each point where its maximal partial sum exceeds $2H$, the maximal sum associated with either $u$ or $v$ exceeds $H$. Thus one of these two real vectors has maximal sum exceeding $H$ on a set of measure at least $\eta/4$. Its norm is at most one, and normalizing it can only increase the maximal sum. Translating the frequencies to positive integers and inserting missing frequencies with coefficient zero gives \eqref{eq:fourier} with $\delta:=\eta/4$.
\end{proof}

The same estimate holds when the frequencies $1,\ldots,m$ are replaced by an arithmetic progression $a+d,\ldots,a+md$, where $a\in\Z$ and $d$ is a positive integer. Indeed,
\begin{equation}\label{eq:affine}
 \sum_{j\leq q}b_j e((a+d\pi(j))x)
 =e(ax)\sum_{j\leq q}b_j e(\pi(j)dx)
\end{equation}
and $|e(ax)|=1$, while $x\mapsto dx$ preserves measure.

Apply Lemma~\ref{lem:two} to the permutation $\pi$ of Lemma~\ref{lem:fourier} and form the system \eqref{eq:construction} with the resulting $N$ and $\sigma$. Fix any permutation $\tau$ of $[N]$. Lemma~\ref{lem:two} supplies a progression $x_1<\cdots<x_m$ such that $x_{\pi(1)},\ldots,x_{\pi(m)}$ occurs as a subsequence of the frequencies of $\phi_{\tau(1)},\ldots,\phi_{\tau(N)}$ on one copy of $\T$. Write $x_j=a+dj$. Give the function with frequency $x_{\pi(j)}$ on this copy the coefficient $b_j$, and give all other functions coefficient zero. On that copy, the maximal absolute partial sum is the maximum absolute value of the sums in \eqref{eq:affine}. By \eqref{eq:fourier}, it exceeds $H$ on a set of measure at least $\delta/2$ in $\T\times\{0,1\}$. This proves the first assertion with $c:=\delta/2$.

\subsection*{The infinite system}

For each $r\geq1$, let $\Phi_r$ be a system of the form \eqref{eq:construction} furnished by the finite construction with $H=4^r$, and denote its defining permutation by $\sigma_r$. Regard these systems as functions of separate coordinates on
\[
 \Omega:=\prod_{r\geq1}(\T\times\{0,1\}).
\]
Equip $\Omega$ with product probability measure $\mu$. Since every function in $\Phi_r$ has mean zero, the union of these finite systems is an ONS, with every function of absolute value one. Enumerate this union as $\{\psi_n\}_{n\geq1}$. The probability space $\Omega$ may be identified modulo null sets with $[0,1]$.

Fix a permutation $\tau$ of $\N$. Let $I_r$ be the smallest interval containing all positions occupied by members of $\Phi_r$ in the sequence $\psi_{\tau(1)},\psi_{\tau(2)},\ldots$. We may choose an infinite set $J\subseteq\N$ so that the intervals $I_r$, $r\in J$, are pairwise disjoint and ordered from left to right by $r$. Indeed, every finite initial segment of this sequence meets only finitely many of the families. After choosing an interval, we can therefore choose a larger $r$ for which $I_r$ lies entirely to its right. For each $r\in J$, apply the finite assertion to the permutation in which the functions of $\Phi_r$ appear. Multiply the resulting unit coefficient vector by $2^{-r}$, and set all other coefficients equal to zero. These coefficients are square-summable, since their squared sum is $\sum_{r\in J}4^{-r}$.

The finite assertion supplies an event $E_r$, depending only on the $r$th coordinate, with $\mu(E_r)\geq\delta/2$. Within $I_r$, only the functions of $\Phi_r$ carry nonzero coefficients, since the intervals $I_{r'}$ with $r'\in J$ are disjoint and all remaining coefficients vanish. On $E_r$, the partial sum at some position inside $I_r$ therefore differs by more than $2^{-r}4^r=2^r$ from the partial sum immediately preceding $I_r$, so at least one partial sum has absolute value greater than $2^{r-1}$. The events $E_r$ are independent because they depend on distinct coordinates, and their measures have divergent sum. The second Borel--Cantelli lemma therefore shows that their limsup has full measure. At every point of this limsup the partial sums are unbounded. By the Riesz--Fischer theorem, the series is the orthogonal expansion of an $L^2$ function. This proves the divergence assertion in the second part of the theorem.

\subsection*{Real-valued systems}

For the real-valued case, adjoin an independent variable $t\in\T$ to each coordinate and replace the finite functions by
\[
 f_n(x,\epsilon,t):=\sqrt2\operatorname{Re}\bigl(e(t)\phi_n(x,\epsilon)\bigr).
\]
Averaging in $t$ shows that these functions are orthonormal, and each is bounded by $\sqrt2$. The coefficients used above are real. If a complex partial sum $S$ has $|S|>H$, then $\sqrt2|\operatorname{Re}(e(t)S)|>H$ for at least half the values of $t$. Thus every permutation of the $f_n$ admits a unit coefficient vector whose maximal sum exceeds $H$ on a set of measure at least $\delta/4$.

The resulting finite systems have mean zero, so the same construction on a product probability space gives the infinite real-valued example. We may take $c:=\delta/4$ throughout.

\begin{remark}[Completeness]
Extend each $\sigma_r$ to a permutation of $\Z$ by fixing $0$ and the unused positive integers and setting $\sigma_r(-n)=-\sigma_r(n)$ for $n\geq1$. Let $n$ range over $\Z$ in \eqref{eq:construction}, write the resulting functions as $u_n$, and adjoin
\[
 v_n(x,\epsilon):=(-1)^\epsilon u_n(x,\epsilon).
\]
Then
\[
 \langle u_n,u_m\rangle=\langle v_n,v_m\rangle=\delta_{nm},
 \qquad
 \langle u_n,v_m\rangle
 =\frac12\bigl(\delta_{nm}-\delta_{\sigma_r(n),\sigma_r(m)}\bigr)=0.
\]
Moreover,
\[
 \frac{u_n+v_n}{2}=\one_{\{\epsilon=0\}}e(nx),
 \qquad
 \frac{u_n-v_n}{2}=\one_{\{\epsilon=1\}}e(\sigma_r(n)x).
\]
Since $\sigma_r$ permutes $\Z$, the functions $u_n$ and $v_n$ form an orthonormal basis on the $r$th coordinate. Their finite tensor products give a complete complex system with absolute value one.

For the real example, first take products of these coordinate basis functions with the characters $e(kt)$, $k\in\Z$, of the additional variable $t$. Form finite tensor products and replace nonreal conjugate pairs by normalized real and imaginary parts. The bound is $\sqrt2$. Setting the additional coefficients to zero retains the divergence property, so the systems in the second assertion may indeed be taken complete.
\end{remark}

\section*{Commentary, acknowledgements, and AI usage}

I would like to thank Jean Bourgain for encouraging me to work on this problem. I first learned of it from his paper \cite{Bourgain} while a graduate student and have been actively thinking about it ever since. When I first met Jean at the Institute for Advanced Study, I asked him about the problem. Although he never explicitly expressed an opinion on whether the answer was positive or negative, our conversations left me with the impression that he was optimistic about a positive answer. He suggested that ideas from the then-recent solution of the Kadison--Singer problem \cite{MSS} might be useful in this direction.

\Needspace{6\baselineskip}
I worked extensively on this project using ChatGPT as a proof assistant and sounding board. Over the course of hundreds of prompts and interactions, I developed a much more complicated counterexample based on Walsh series and the hypergraph container results of Saxton and Thomason \cite{SaxtonThomason}. After further reflection on that example, I found the approach presented here, which is much simpler and more intuitive.

\begingroup
\small
\begin{thebibliography}{99}
\setlength{\itemsep}{0pt}
\bibitem{Billard}
P. Billard, \emph{Sur la convergence presque partout des s\'eries de Fourier-Walsh des fonctions de l'espace $L^2(0,1)$}, Studia Math. \textbf{28} (1967), 363--388.
\href{https://doi.org/10.4064/sm-28-3-363-388}{doi:10.4064/sm-28-3-363-388}.

\bibitem{BKL}
R. G. Bilyeu, R. R. Kallman and P. W. Lewis,
\emph{Rearrangements and category}, Pacific J. Math. \textbf{121} (1986), no.~1, 41--46.
\href{https://msp.org/pjm/1986/121-1/pjm-v121-n1-p04-s.pdf}{Article}.

\bibitem{Bourgain}
J. Bourgain, \emph{On Kolmogorov's rearrangement problem for orthogonal systems and Garsia's conjecture}, in Geometric Aspects of Functional Analysis (1987--88), Lecture Notes in Math. \textbf{1376}, Springer, Berlin, 1989, 209--250.
\href{https://doi.org/10.1007/BFb0090057}{doi:10.1007/BFb0090057}.

\bibitem{BourgainLambda}
J. Bourgain, \emph{Bounded orthogonal systems and the $\Lambda(p)$-set problem}, Acta Math. \textbf{162} (1989), 227--245.
\href{https://doi.org/10.1007/BF02392838}{doi:10.1007/BF02392838}.

\bibitem{Carleson}
L. Carleson, \emph{On convergence and growth of partial sums of Fourier series}, Acta Math. \textbf{116} (1966), 135--157.
\href{https://doi.org/10.1007/BF02392815}{doi:10.1007/BF02392815}.

\bibitem{Garsia1964}
A. M. Garsia, \emph{Existence of almost everywhere convergent rearrangements for Fourier series of $L^2$ functions}, Ann. of Math. (2) \textbf{79} (1964), 623--629.
\href{https://doi.org/10.2307/1970412}{doi:10.2307/1970412}.

\bibitem{Garsia1970}
A. M. Garsia, \emph{Topics in Almost Everywhere Convergence}, Lectures in Advanced Mathematics \textbf{4}, Markham, Chicago, 1970.

\bibitem{Karagulyan}
G. A. Karagulyan, \emph{On Weyl multipliers of the rearranged trigonometric system}, Mat. Sb. \textbf{211} (2020), no.~12, 49--82; English transl., Sb. Math. \textbf{211} (2020), 1704--1736.
\href{https://arxiv.org/abs/2004.01003}{arXiv:2004.01003}.

\bibitem{Kolmogorov1924}
A. Kolmogoroff, \emph{Une contribution \`a l'\'etude de la convergence des s\'eries de Fourier}, Fund. Math. \textbf{5} (1924), 96--97.
\href{https://matwbn.icm.edu.pl/ksiazki/fm/fm5/fm5111.pdf}{Article}.

\bibitem{KolmogorovMenshov}
A. Kolmogoroff and D. Menchoff, \emph{Sur la convergence des s\'eries de fonctions orthogonales}, Math. Z. \textbf{26} (1927), 432--441.
\href{https://doi.org/10.1007/BF01475463}{doi:10.1007/BF01475463}.

\bibitem{MSS}
A. W. Marcus, D. A. Spielman and N. Srivastava, \emph{Interlacing families II: Mixed characteristic polynomials and the Kadison--Singer problem}, Ann. of Math. (2) \textbf{182} (2015), no.~1, 327--350.
\href{https://doi.org/10.4007/annals.2015.182.1.8}{doi:10.4007/annals.2015.182.1.8}.

\bibitem{Olevskii1964}
A. M. Olevskii, \emph{On complete systems of convergence}, Dokl. Akad. Nauk SSSR \textbf{159} (1964), no.~1, 21--24; English transl., Soviet Math. Dokl. \textbf{5} (1964), 1425--1429.
\href{https://www.mathnet.ru/eng/dan30299}{Article}.

\bibitem{OlevskiiBook}
A. M. Olevskii, \emph{Fourier Series with Respect to General Orthogonal Systems}, Ergebnisse der Mathematik und ihrer Grenzgebiete \textbf{86}, Springer-Verlag, Berlin--Heidelberg--New York, 1975.
\href{https://link.springer.com/book/9783642660580}{Book}.

\bibitem{SaxtonThomason}
D. Saxton and A. Thomason, \emph{Hypergraph containers}, Invent. Math. \textbf{201} (2015), no.~3, 925--992.
\href{https://arxiv.org/abs/1204.6595}{arXiv:1204.6595}.

\bibitem{Szemeredi}
E. Szemer\'edi, \emph{On sets of integers containing no $k$ elements in arithmetic progression}, Acta Arith. \textbf{27} (1975), 199--245.
\href{https://doi.org/10.4064/aa-27-1-199-245}{doi:10.4064/aa-27-1-199-245}.

\bibitem{Zahorski}
Z. Zahorski, \emph{Une s\'erie de Fourier permut\'ee d'une fonction de classe $L^2$ divergente presque partout}, C. R. Acad. Sci. Paris \textbf{251} (1960), 501--503.

\end{thebibliography}

\medskip
\noindent\textsc{Enfield, New Hampshire, USA}\\[2pt]
\textit{Email address:} \href{mailto:mlewko@gmail.com}{\texttt{mlewko@gmail.com}}
\endgroup
\end{document}
